\documentclass[11pt,a4paper]{article}
\usepackage[margin=2.5cm]{geometry}
\usepackage{amsmath,amssymb,amsthm}
\usepackage{booktabs}
\usepackage{graphicx}
\usepackage{hyperref}
\usepackage{float}
\usepackage{url}

\newcommand{\rad}{\operatorname{rad}}
\newcommand{\Jac}{\operatorname{Jac}}
\newcommand{\Cond}{\operatorname{Cond}}
\newcommand{\End}{\operatorname{End}}
\newcommand{\Res}{\operatorname{Res}}
\newcommand{\sw}{\operatorname{sw}}
\newcommand{\Gm}{\mathbb{G}_{\mathrm{m}}}
\newcommand{\Ga}{\mathbb{G}_{\mathrm{a}}}

\newtheorem{theorem}{Theorem}[section]
\newtheorem{proposition}[theorem]{Proposition}
\newtheorem{corollary}[theorem]{Corollary}
\newtheorem{conjecture}[theorem]{Conjecture}
\newtheorem{lemma}[theorem]{Lemma}
\theoremstyle{definition}
\newtheorem{definition}[theorem]{Definition}
\theoremstyle{remark}
\newtheorem{remark}[theorem]{Remark}

\title{The $\gamma$-Invariant of Goldbach--Frey Curves:\\
  Cluster Depths and Conjectural Bounds\\
  for the Wild Conductor at $2$}
\author{Ruqing Chen\\[4pt]
  \small GUT Geoservice Inc., Montreal\\
  \small\texttt{ruqing@hotmail.com}}
\date{February 2026}

\begin{document}
\maketitle

\begin{abstract}
For the Goldbach--Frey curve
$C\colon y^2 = x(x^2 - p^2)(x^2 - q^2)$ with $p \ne q$
distinct odd primes, the conductor exponent $f_2$ at the
prime~$2$ takes values in $\{4, 7, 8\}$ and has resisted
closed-form description.  We prove a structural lemma
showing that exactly one of $v_2(p-q)$ and $v_2(p+q)$
equals~$1$, and define
$\gamma = \max(v_2(p-q),\, v_2(p+q)) \ge 2$.
We conjecture that $\gamma$ completely determines $f_2$:
$f_2 = 8$ if $\gamma = 2$, $f_2 = 7$ if $\gamma = 3$,
and $f_2 = 4$ if $\gamma \ge 4$.
The case $\gamma \ge 4$ is \emph{proved}: the Weil
restriction $\Jac(C) \sim E_1 \times E_2$
(Paper~\cite{Chen2026WR}) gives
$f_2 = f_2(E_1) + f_2(E_2)$, and a $j$-invariant
analysis shows $v_2(j) = 6 - 2\gamma < 0$, placing
both curves in the potentially multiplicative regime
where $f_2(E) = 2$ unconditionally.
The cases $\gamma \le 3$ are verified against all $10$
Magma-computed conductors but remain conjectural.
We note, however, that since
$f_2(\Jac) = f_2(E_1) + f_2(E_2)$, the $10$ Magma
values can in principle be independently verified
by applying Tate's algorithm to the elliptic
curves~$E_1$ and~$E_2$ over~$\mathbb{Q}_2$---an
unconditional computation with no correctness caveats,
unlike Magma's genus-$2$ Ogg heuristic.
The $2$-adic cluster tree of the Weierstrass roots
provides geometric intuition for the formula, although
we emphasise that the Dokchitser cluster picture theory
does not directly compute the Swan conductor at $p = 2$
(where wild ramification requires Artin--Schreier
analysis beyond root distances alone).
We also note that the Magma values for $\gamma \le 3$
rely on the Ogg formula heuristic at $v_2(\Delta) \ge 12$,
so the $\gamma$-conjecture at $\gamma \le 3$ should be
regarded as empirically robust but not yet rigorously
established.
A proof for the remaining cases requires computing
the Dokchitser defect over~$\mathbb{Z}_2$, which we
pose as an open problem.
\end{abstract}


% ====================================================================
\section{Introduction}

Papers~\cite{Chen2026TCV,Chen2026UTS} established that
the conductor of $\Jac(C)$ at every odd prime~$r$
dividing the discriminant is $f_r = 2$, yielding the
odd-part formula
$\Cond_{\mathrm{odd}} =
[\rad_{\mathrm{odd}}(pqN|p{-}q|)]^2$ where $N = (p+q)/2$.
The conductor exponent $f_2$ at the prime~$2$, however,
remained a black box: Magma computations on $10$ test
curves revealed $f_2 \in \{4, 7, 8\}$ with no apparent
pattern.

The difficulty at $r = 2$ is fundamental: the curve~$C$
has wild ramification, so the tame conductor formula no
longer suffices.  The cluster picture
of Dokchitser~\cite{Dokchitser} computes the
conductor of hyperelliptic curves from the $p$-adic
distances between the Weierstrass roots, and provides
geometric intuition for our formula.  However, we
emphasise two important caveats:

\begin{enumerate}
\item[(a)] The cluster picture at $p = 2$ does not
  directly determine the Swan conductor.  Since the
  hyperelliptic degree~$2$ coincides with the residue
  characteristic, the wild ramification requires
  Artin--Schreier analysis and the computation of
  the \emph{Dokchitser defect}---a correction term
  that depends on higher-order congruences of the
  polynomial coefficients, not merely on root distances.
\item[(b)] The Magma values $f_2 \in \{4, 7, 8\}$
  themselves rely on the Ogg--Saito formula, which
  carries an internal warning when
  $v_2(\Delta) \ge 12$: ``no correctness guarantee.''
  Since $v_2(\Delta)$ is itself determined by~$\gamma$,
  the $\gamma \to f_2$ correspondence may partially
  reflect the internal logic of Magma's heuristic
  rather than the true geometric conductor.
\end{enumerate}

\noindent
With these caveats in mind, our results are:
\begin{enumerate}
\item prove a structural lemma on $v_2(p \pm q)$
  (Lemma~\ref{lem:structure});
\item conjecture a closed-form formula for $f_2$
  in terms of
  $\gamma = \max(v_2(p{-}q), v_2(p{+}q))$
  (Conjecture~\ref{conj:f2}), verified on all
  $10$ Magma data points;
\item explain the formula via the $2$-adic cluster tree
  (Section~\ref{sec:cluster});
\item establish the density of each conductor value
  (Proposition~\ref{prop:density});
\item state the complete conductor formula combining
  all papers (Section~\ref{sec:complete}).
\end{enumerate}


% ====================================================================
\section{The Structural Lemma}\label{sec:lemma}

\begin{lemma}\label{lem:structure}
  For distinct odd primes $p \ne q$, exactly one of
  $v_2(p - q)$ and $v_2(p + q)$ equals~$1$, and the
  other is $\ge 2$.
\end{lemma}

\begin{proof}
  Write $p = 2a + 1$ and $q = 2b + 1$.  Then
  $p - q = 2(a - b)$ and $p + q = 2(a + b + 1)$.
  Since $(a - b) + (a + b + 1) = 2a + 1$ is odd,
  the integers $a - b$ and $a + b + 1$ have different
  parity.  Exactly one is odd (contributing $v_2 = 1$
  to the corresponding $p \pm q$) and the other is
  even (contributing $v_2 \ge 2$).
\end{proof}

\begin{definition}\label{def:gamma}
  The \emph{cluster depth parameter} is
  $\gamma = \max(v_2(p - q),\, v_2(p + q)) \ge 2$.
\end{definition}


% ====================================================================
\section{The Wild Conductor Formula}\label{sec:formula}

\begin{conjecture}\label{conj:f2}
  For the Goldbach--Frey curve~$C$ with $p \ne q$
  distinct odd primes,
  \begin{equation}\label{eq:f2}
    f_2 \;=\; \begin{cases}
      8 & \text{if } \gamma = 2, \\
      7 & \text{if } \gamma = 3, \\
      4 & \text{if } \gamma \ge 4.
    \end{cases}
  \end{equation}
\end{conjecture}

\noindent
The case $\gamma \ge 4$ is proved unconditionally in
Proposition~\ref{prop:stable} below; only the cases
$\gamma = 2$ and $\gamma = 3$ remain conjectural.

\begin{table}[H]
\centering
\begin{tabular}{rrrrrrcl}
\toprule
$p$ & $q$ & $v_2(p{-}q)$ & $v_2(p{+}q)$ & $\gamma$
& $f_2^{\mathrm{Magma}}$ & $f_2^{\mathrm{pred}}$ & \\
\midrule
\multicolumn{8}{l}{\small\itshape Conjectural regime
  ($\gamma \le 3$, potentially good reduction):} \\[2pt]
 3 &  7 & 2 & 1 & 2 & 8 & 8 & $\checkmark$ \\
13 & 17 & 2 & 1 & 2 & 8 & 8 & $\checkmark$ \\
 3 & 17 & 1 & 2 & 2 & 8 & 8 & $\checkmark$ \\
 7 & 19 & 2 & 1 & 2 & 8 & 8 & $\checkmark$ \\
 7 & 53 & 1 & 2 & 2 & 8 & 8 & $\checkmark$ \\
11 & 19 & 3 & 1 & 3 & 7 & 7 & $\checkmark$ \\
 3 & 37 & 1 & 3 & 3 & 7 & 7 & $\checkmark$ \\
\midrule
\multicolumn{8}{l}{\small\itshape Proved regime
  ($\gamma \ge 4$, Proposition~\ref{prop:stable}):} \\[2pt]
 7 & 23 & 4 & 1 & 4 & 4 & 4 & $\checkmark$ \\
 7 & 41 & 1 & 4 & 4 & 4 & 4 & $\checkmark$ \\
11 & 37 & 1 & 4 & 4 & 4 & 4 & $\checkmark$ \\
\bottomrule
\end{tabular}
\caption{Verification of Conjecture~\ref{conj:f2}
  against all $10$ Magma-computed conductors.}
\label{tab:verification}
\end{table}

\begin{proposition}[Stabilisation at $\gamma \ge 4$]
\label{prop:stable}
  For all $\gamma \ge 4$, the conductor exponent
  $f_2 = 4$.
\end{proposition}

\begin{proof}
  By Paper~\cite{Chen2026WR}, $\Jac(C) \sim E_1 \times E_2$
  over~$\mathbb{Q}$, where
  $E_1\colon Y^2 = X(X - p^2)(X - q^2)$ and
  $E_2 = E_1^{(-1)}$.  Since
  $V_\ell(\Jac) \cong V_\ell(E_1) \oplus V_\ell(E_2)$
  as $G_\mathbb{Q}$-representations, the conductor
  exponent is additive:
  $f_2(\Jac) = f_2(E_1) + f_2(E_2)$.

  For both $E_1$ and $E_2$, we compute:
  $c_4 = 16(p^4 - p^2 q^2 + q^4)$ with
  $v_2(c_4) = 4$ (since $p^4 - p^2 q^2 + q^4 \equiv 1
  \pmod{8}$ for all odd $p, q$), and
  $\Delta = 16 p^4 q^4 (p^2 - q^2)^2$ with
  $v_2(\Delta) = 6 + 2\gamma$.
  The $j$-invariant satisfies
  $v_2(j) = 3 v_2(c_4) - v_2(\Delta) = 6 - 2\gamma$.

  For $\gamma \ge 4$, $v_2(j) \le -2 < 0$, so
  both $E_1$ and~$E_2$ have \emph{potentially
  multiplicative} reduction at~$2$ (Kodaira type
  $I_n^*$ for some $n \ge 0$).  By Ogg's theorem,
  the conductor exponent of any elliptic curve with
  Kodaira type $I_n^*$ is $f_p = 2$ at every
  prime~$p$, including $p = 2$ (the wild part
  vanishes for this reduction type).
  Therefore $f_2(\Jac) = 2 + 2 = 4$ for all
  $\gamma \ge 4$.
\end{proof}

\begin{remark}[The odd value $f_2 = 7$]
\label{rem:odd}
  Since $f_2(\Jac) = f_2(E_1) + f_2(E_2)$ and
  $f_2 = 7$ is odd, the decomposition at $\gamma = 3$
  must be asymmetric: $f_2(E_1) \ne f_2(E_2)$ (e.g.\
  $4 + 3$ or $3 + 4$).  This is explained by the fact
  that $E_2 = E_1^{(-1)}$ is the quadratic twist by~$-1$,
  and $\mathbb{Q}_2(\sqrt{-1})/\mathbb{Q}_2$ is a
  \emph{ramified} quadratic extension (discriminant~$-4$).
  Twisting an elliptic curve with additive reduction by
  a ramified character routinely alters the Swan
  conductor~\cite{Kramer1981}.
  At $\gamma = 2$, $f_2 = 8$ is even,
  consistent with $f_2(E_1) = f_2(E_2) = 4$ (symmetric).
\end{remark}

\begin{remark}[$j$-invariant regime transition]
\label{rem:regime}
  The formula $v_2(j) = 6 - 2\gamma$ reveals a sharp
  phase transition at $\gamma = 3 \leftrightarrow 4$:
  for $\gamma \le 3$, $E_1$ and~$E_2$ have potentially
  good reduction ($v_2(j) \ge 0$), where wild
  ramification can occur; for $\gamma \ge 4$,
  the reduction is potentially multiplicative
  ($v_2(j) < 0$), where the Swan conductor vanishes.
  The large jump $f_2 = 7 \to 4$ at this threshold
  reflects the disappearance of all wild ramification.
\end{remark}


% ====================================================================
\section{The Conductor Decomposition}\label{sec:decomp}

The conductor exponent at a prime~$r$ decomposes as
\begin{equation}\label{eq:decomp}
  f_r \;=\; \underbrace{\dim V_\ell
    - \dim V_\ell^{I_r}}_{\text{tame part}}
  \;+\; \underbrace{\sw(V_\ell \mid_{I_r})}_{\text{wild part}},
\end{equation}
where $V_\ell = T_\ell(\Jac(C)) \otimes_{\mathbb{Z}_\ell}
\mathbb{Q}_\ell$ is the $\ell$-adic Tate module
($\ell \ne r$),
$I_r$ is the inertia group, and $\sw$ is the Swan
conductor.

For $r = 2$, the Magma data is consistent with the
hypothesis that
$V_\ell^{I_2} = 0$ (no inertia invariants), which would
give a tame part of~$4$ and yield:
\begin{equation}\label{eq:swan}
  f_2 = 4 + \sw, \qquad
  \sw \in \{0, 3, 4\}.
\end{equation}

\begin{remark}
  The vanishing $V_\ell^{I_2} = 0$ has not been proved
  in this paper; it is an assumption consistent with
  $f_2 \ge 4$ in all observed cases.  A proof would
  require an explicit analysis of the N\'eron model
  of $\Jac(C)$ over~$\mathbb{Z}_2$, which is
  substantially harder than the tame case treated
  in Paper~\cite{Chen2026UTS}.
\end{remark}

\noindent
If the decomposition $f_2 = 4 + \sw$ holds, the three
observed values correspond to:
$\sw = 0$ ($\gamma \ge 4$, $f_2 = 4$):
vanishing Swan conductor;
$\sw = 3$ ($\gamma = 3$, $f_2 = 7$);
$\sw = 4$ ($\gamma = 2$, $f_2 = 8$):
maximal wild part for this family.


% ====================================================================
\section{The Cluster Picture at $r = 2$}\label{sec:cluster}

\subsection{The $2$-adic distance matrix}

The Weierstrass roots of $f(x) = x(x^2 - p^2)(x^2 - q^2)$
are $\mathcal{R} = \{0, p, -p, q, -q\}$, all $2$-adic
units except~$0$.  The pairwise distances are
$v_2(p - (-p)) = 1$, $v_2(q - (-q)) = 1$, and by
Lemma~\ref{lem:structure}, exactly one of
$v_2(p{-}q), v_2(p{+}q)$ equals~$\gamma$ and the other
equals~$1$.

\subsection{The cluster tree}

WLOG assume $v_2(p - q) = \gamma$ (the case
$v_2(p + q) = \gamma$ is symmetric).  The cluster tree
over~$\mathbb{Q}_2$ has three levels:

\textbf{Level 0:}
  Root $0$ is isolated.  The four roots
  $\{p, -p, q, -q\}$ form a cluster~$\mathcal{C}_0$
  of depth~$1$.

\textbf{Level 1:}
  Within $\mathcal{C}_0$, the pairs $\{p, q\}$ and
  $\{-p, -q\}$ form sub-clusters (each of depth~$\gamma$).

\textbf{Level $\gamma$:}
  Each sub-cluster splits into singletons.

\begin{figure}[H]
  \centering
  \includegraphics[width=\textwidth]{../figures/fig_cluster_tree.pdf}
  \caption{(a)~Cluster trees for $\gamma = 2, 3, \ge 4$.
    (b)~Distribution of $\gamma$ matching
    $P(\gamma = k) \approx 2^{-k}$.}
  \label{fig:cluster}
\end{figure}

\subsection{From cluster depth to wild conductor:
  heuristic interpretation}

The cluster tree provides a compelling geometric
\emph{heuristic} for the $\gamma$-conjecture, but we
stress that it does not constitute a proof.

At odd primes, the cluster picture directly computes
the conductor via the formula of
Dokchitser~\cite{Dokchitser}.  At $p = 2$, however,
the wild ramification introduces a \emph{defect} term
that depends on higher-order congruence properties of
the polynomial coefficients (specifically, their
reduction modulo $4$, $8$, \ldots), not merely on the
$2$-adic distances between roots.  Computing this defect
for our family requires an explicit Artin--Schreier
analysis or a full stable model computation
over~$\mathbb{Z}_2$ via the Obus--Wewers
theory~\cite{ObusWewers}.

Nevertheless, the qualitative principle is suggestive:
deeper clusters $\Longrightarrow$ roots closer
$2$-adically $\Longrightarrow$ less wild ramification
$\Longrightarrow$ smaller~$f_2$.
The quantitative jumps $8 \to 7 \to 4$ at
$\gamma = 2, 3, \ge 4$ remain conjectural.

\begin{remark}[Magma's internal algorithm]\label{rem:magma}
  The Magma values $f_2 \in \{4, 7, 8\}$ that we use as
  ground truth are themselves computed via the Ogg--Saito
  formula, which Magma flags with a warning when
  $v_2(\Delta) \ge 12$.  Since $v_2(\Delta)$ is a
  deterministic function of~$\gamma$ for our family
  ($v_2(\Delta) = 6 + 2\gamma$ for the elliptic
  quotients, by Proposition~\ref{prop:stable}), the
  perfect $\gamma \to f_2$ correspondence may partially
  reflect the internal logic of Magma's Ogg heuristic.
  Independent verification (e.g., via \'etale cohomology
  computations in SageMath or direct $\ell$-adic Galois
  image calculations) would strengthen the conjecture.
\end{remark}


% ====================================================================
\section{Density of Conductor Values}\label{sec:density}

\begin{proposition}\label{prop:density}
  Among Goldbach pairs $(p, q)$, the asymptotic density is
  $P(f_2 = 8) = 1/2$, $P(f_2 = 7) = 1/4$,
  $P(f_2 = 4) = 1/4$.
\end{proposition}

\begin{proof}
  By Lemma~\ref{lem:structure}, $\gamma$ equals the
  $2$-adic valuation of whichever of $p \pm q$ is more
  divisible by~$2$.
  By the Prime Number Theorem in arithmetic progressions,
  the residue classes of primes $p$ and~$q$ modulo~$2^k$
  are asymptotically uniformly distributed among the odd
  classes for any fixed~$k$.  The $2$-adic valuations of
  $p \pm q$ therefore behave as those of a random even
  integer, giving
  $P(\gamma = 2) = 1/2$, $P(\gamma = 3) = 1/4$,
  $P(\gamma \ge 4) = 1/4$.
\end{proof}

\noindent
Verified on $425{,}082$ pairs: $50.0\%$ / $25.0\%$ / $25.0\%$.

\begin{figure}[H]
  \centering
  \includegraphics[width=\textwidth]{../figures/fig_wild_conductor.pdf}
  \caption{(a)~$f_2$ as a step function of $\gamma$.
    (b)~Distribution: exactly $50\%$/$25\%$/$25\%$.}
  \label{fig:wild}
\end{figure}


% ====================================================================
\section{The Complete Conductor Formula}\label{sec:complete}

\begin{conjecture}[Complete conductor formula]
\label{conj:complete}
  For the Goldbach--Frey curve with $p + q = 2N$:
  \begin{equation}\label{eq:complete}
    \boxed{N_A \;=\; 2^{f_2(\gamma)} \cdot
    \bigl[\rad_{\mathrm{odd}}(p \cdot q \cdot N
    \cdot |p - q|)\bigr]^2}
  \end{equation}
  where $\gamma = \max(v_2(p{-}q), v_2(p{+}q))$ and
  $f_2(\gamma)$ is given by~\eqref{eq:f2}.
\end{conjecture}

\noindent
The $\gamma \ge 4$ regime is unconditionally proved
(Proposition~\ref{prop:stable}).  Subject to the
$\gamma$-conjecture for $\gamma \le 3$, every ingredient
is explicit:
$f_r = 2$ at odd primes (Paper~\cite{Chen2026UTS}, proved);
$\rad_{\mathrm{odd}}$ (Paper~\cite{Chen2026TCV}, verified
on $425{,}082$ pairs); $f_2(\gamma)$ (this paper:
proved for $\gamma \ge 4$, conjectural for $\gamma \le 3$,
$10/10$ against Magma).


% ====================================================================
\section{Discussion}

The wild conductor at $r = 2$ was the last remaining
unknown in the arithmetic of Goldbach--Frey curves.
The $\gamma$-invariant
$\gamma = \max(v_2(p{-}q), v_2(p{+}q))$ provides
a clean conjectural description: the three-valued
function $f_2 \in \{4, 7, 8\}$ appears to be
controlled by a single $2$-adic parameter measuring
how close $p$ and~$q$ are in the $2$-adic metric.

The qualitative pattern is:
\[
  \gamma \uparrow
  \;\Longrightarrow\;
  \text{deeper cluster}
  \;\Longrightarrow\;
  f_2 \downarrow.
\]
The conjectured threshold at $\gamma = 4$, where $f_2$
stabilises at its minimum value~$4$, is now
\emph{proved} (Proposition~\ref{prop:stable}):
the $j$-invariant satisfies $v_2(j) = 6 - 2\gamma$,
so for $\gamma \ge 4$ both $E_1$ and~$E_2$ have
potentially multiplicative reduction, forcing
$f_2(E) = 2$ per curve and $f_2(\Jac) = 4$.

Two obstacles remain for a complete proof of the
$\gamma$-conjecture:
\begin{enumerate}
\item For $\gamma \le 3$ (potentially good reduction,
  $v_2(j) \ge 0$), the cluster picture does not
  directly yield the Swan conductor without computing
  the Dokchitser defect, which requires
  Artin--Schreier analysis beyond root distances.
\item The Magma ground truth for $\gamma \le 3$ relies
  on the Ogg heuristic with a correctness caveat at
  high $v_2(\Delta)$ (Remark~\ref{rem:magma}).
\end{enumerate}
Note that the $\gamma \ge 4$ regime is now
unconditional---it does not depend on Magma data at all.
A complete proof for $\gamma \le 3$ would require
either a full stable model computation
over~$\mathbb{Z}_2$ using the Obus--Wewers
theory~\cite{ObusWewers}, or independent verification
of $f_2$ via $\ell$-adic Galois image computations.
We pose this as an open problem.


% ====================================================================
\section*{Acknowledgments}

Scripts and data are at
\url{https://github.com/Ruqing1963/goldbach-wild-conductor}.


\begin{thebibliography}{99}

\bibitem{Chen2026TCV}
R.~Chen,
\emph{The true conductor of Goldbach--Frey curves},
Zenodo, 2026.  \url{https://zenodo.org/records/18749731}

\bibitem{Chen2026UTS}
R.~Chen,
\emph{Universal tame semistability of Goldbach--Frey Jacobians},
Zenodo, 2026.  \url{https://zenodo.org/records/18751169}

\bibitem{Chen2026Census}
R.~Chen,
\emph{A conductor census of 425,082 Goldbach--Frey curves},
Zenodo, 2026.  \url{https://zenodo.org/records/18751442}

\bibitem{Chen2026WR}
R.~Chen,
\emph{Goldbach--Frey Jacobians as Weil restrictions},
Zenodo, 2026.  \url{https://zenodo.org/records/18752179}

\bibitem{Dokchitser}
T.~Dokchitser,
Models of hyperelliptic curves,
\emph{Trans.\ Amer.\ Math.\ Soc.}~\textbf{372} (2019),
4519--4577.

\bibitem{DD2022}
T.~Dokchitser and V.~Dokchitser,
Quotients of hyperelliptic curves and \'etale cohomology,
\emph{Q. J. Math.}~\textbf{73} (2022), 747--768.

\bibitem{ObusWewers}
A.~Obus and S.~Wewers,
Explicit resolution of weak wild quotient singularities,
\emph{J.\ Algebraic Geom.}~\textbf{29} (2020), 691--728.

\bibitem{Kramer1981}
K.~Kramer,
Arithmetic of elliptic curves upon quadratic extension,
\emph{Trans.\ Amer.\ Math.\ Soc.}~\textbf{264} (1981),
121--135.

\end{thebibliography}

\end{document}
